\chapter{Side-Effect Monad Ordering}
\label{chap:side_effects_monads}

\section{Reconciling Pure Functional Dataflow with I/O}
Pure dataflow languages traditionally struggle with imperative I/O operations. In a pure dataflow graph (such as Sisal 1.2), node execution is driven strictly by data dependencies: a node executes as soon as all of its input arguments become available. When sub-expressions have no data dependencies between them, they are scheduled concurrently across available processor threads or executed in arbitrary order.

While this non-deterministic out-of-order execution model is ideal for high-performance data parallelism, it presents severe challenges for practical programming tasks:
\begin{itemize}
    \item \textbf{Console Output Interleaving}: Concurrent printing from multiple parallel sub-graph branches causes character streams to collide and interleave unpredictably.
    \item \textbf{Debug Tracing}: Programmers cannot easily trace algorithm iteration steps or log intermediate values when execution timing varies between runs.
    \item \textbf{Sequential File/Stream Access}: Operations like writing consecutive record chunks to stdout or disk files require a guaranteed, deterministic order.
\end{itemize}

Sisal-2026 solves this challenge by introducing \textbf{Monad Control Ports} into the compiler's IF1 intermediate representation graph.

\section{What is a Monad in Sisal-2026?}
In functional programming and category theory, a \emph{monad} is a design pattern that structures computations by encapsulating value evaluation alongside side-effect execution contexts. In Sisal-2026, a monad represents an implicit, linearly-threaded sequencing state that enforces a strict execution order between side-effecting operations without violating the language's core single-assignment immutability rules.

Conceptually, a side-effecting built-in function in Sisal-2026 operates as a state-transformer monad:
\begin{equation}
f :: A \to S \to (B, S)
\end{equation}
where $A$ is the input argument, $B$ is the returned result value, and $S$ represents an abstract "World State" token. Rather than exposing state tokens to the programmer, the compiler automatically threads this state token through side-effecting operations, creating an explicit linear execution sequence.

\section{Serialization \& Imperative Style for Practical Scenarios}
To handle practical software engineering requirements—such as progress logging, error reporting, and stream output—Sisal-2026 enables developers to write familiar, imperative-style sequential code within pure functional routines.

When a developer writes consecutive side-effecting statements:
\begin{lstlisting}[caption={Imperative-Style Logging Sequence in Sisal-2026}]
function ExecutePipeline( data : array_dv[double] returns array_dv[double] )
  let
    _ := printf("Phase 1: Starting computation...\n");
    R1 := ProcessPhase1(data);
    _ := printf("Phase 2: Intermediate processing complete.\n");
    R2 := ProcessPhase2(R1);
    _ := printf("Phase 3: Finalizing results.\n")
  in
    R2
  end let
end function
\end{lstlisting}

Sisal-2026 reconciles this imperative notation with its dataflow engine through \textbf{deterministic serialization}:
\begin{enumerate}
    \item \textbf{Localized Imperative Sequencing}: Side-effecting statements are executed in exact lexicographic (source file) order.
    \item \textbf{Concurrent Functional Subgraphs}: Pure mathematical expressions (\texttt{ProcessPhase1}, \texttt{ProcessPhase2}, array arithmetic) continue to execute with maximum data parallelism across all available CPU worker threads.
    \item \textbf{No Non-Deterministic Race Conditions}: Output buffers are never corrupted by out-of-order thread scheduling.
\end{enumerate}

\section{Monad Control Ports in the IF1 Intermediate Representation}
A \textbf{Monad Control Port} is a compiler-inserted IR graph port designed to carry sequence control tokens. When the Sisal-2026 parser encounters side-effecting built-in procedures (\texttt{printf}, \texttt{cout}, \texttt{cerr}), it decorates the corresponding IF1 AST nodes with dedicated monad ports:
\begin{itemize}
    \item \texttt{PRINTF\_TY}: Monad sequence control port for standard output logging.
    \item \texttt{COUT\_TY}: C++ stream output monad sequence port.
    \item \texttt{CERR\_TY}: Standard error stream monad sequence port.
\end{itemize}

During graph lowering, the compiler constructs artificial control-dependency edges (\emph{Monad Control Edges}) between consecutive side-effecting nodes based on their program text position. An IF1 node possessing an incoming Monad Control Edge is prohibited from firing until the predecessor side-effecting node has completed execution and emitted its control token.

\begin{figure}[h!]
\centering
\begin{tikzpicture}[node distance=2.2cm, auto, >=stealth']
    % Nodes
    \node [draw=sisalnavy, rounded corners=4pt, rectangle, fill=sisalblue!15, minimum width=3.2cm, minimum height=1.0cm, align=center, font=\bfseries] (p1) {\texttt{printf("Start")}};
    \node [draw=sisalnavy, rounded corners=4pt, rectangle, fill=sisalgreen!15, minimum width=3.2cm, minimum height=1.0cm, align=center, font=\bfseries, right=3.5cm of p1] (compute) {\texttt{HeavyCompute()}};
    \node [draw=sisalnavy, rounded corners=4pt, rectangle, fill=sisalblue!15, minimum width=3.2cm, minimum height=1.0cm, align=center, font=\bfseries, below=2.2cm of compute] (p2) {\texttt{printf("End")}};
    
    % Dataflow Edges (Pure Computation Path)
    \draw [->, thick, sisaldarkblue] (p1) -- node[above, font=\small\slshape] {Dataflow Edge} (compute);
    \draw [->, thick, sisaldarkblue] (compute) -- node[right, font=\small\slshape] {Data Depend} (p2);
    
    % Monad Sequencing Control Edge (Curved to prevent box collision)
    \draw [->, thick, sisalorange!90!black, dashed] (p1) to[bend right=35] node[below left=2pt, font=\small\bfseries, align=center] {\texttt{PRINTF\_TY}\\[1pt]\scriptsize Monad Control Edge} (p2);
\end{tikzpicture}
\caption{Monad Control Port Dependency Edge in IF1 Graph.}
\label{fig:monad_port}
\end{figure}

As illustrated in Figure~\ref{fig:monad_port}, pure computation paths (\texttt{HeavyCompute()}) execute asynchronously according to data availability, while side-effecting nodes are linked via the dashed \texttt{PRINTF\_TY} Monad Control Edge to guarantee strict ordering.

\section{Value Pass-Through vs. Wildcard Idioms}
Sisal-2026 supports two primary statement idioms for \texttt{printf} logging:
\begin{itemize}
    \item \textbf{Wildcard Binding (\texttt{\_ := printf(...)})}: Evaluates the formatted logging statement and discards the return value using the wildcard pattern \texttt{\_}. This pattern is ideal for stand-alone progress messages.
    \item \textbf{Value Pass-Through (\texttt{res := printf("\%d\textbackslash n", val)})}: Logs the formatted message while transparently evaluating and returning the printed value argument. This permits inline debug logging without disrupting dataflow or polluting the program with temporary variable declarations.
\end{itemize}

\section{Deterministic Lexicographic Ordering Guarantees}
By combining Monad Control Ports in the IR with single-assignment language semantics, Sisal-2026 achieves a powerful synthesis:
\begin{enumerate}
    \item \textbf{Guaranteed Determinism}: Given identical inputs, a Sisal-2026 program will produce the exact same console log output in the exact same sequence across every run, regardless of hardware core count or thread scheduling order.
    \item \textbf{Zero Impact on Pure Parallelism}: Monad control edges sequence only nodes decorated with side-effect ports. Independent mathematical operations, array contractions, tensor operations, and loop iterations remain 100\% parallelizable across multicore targets.
\end{enumerate}
